\clearpage
\begin{examplebox}
\textbf{\textcolor{CExample}{例1（最基本：直接代入）}}\par
\textbf{题目：}抛物线 $y^2=4x$ 的过焦点直线倾角为 $30^\circ$，求焦点弦长。\par
\textbf{解：}由 $2p=4$ 得 $p=2$，于是
\[
\boxed{L=\frac{4}{\sin^230^\circ}=\frac4{1/4}=16.}
\]
\textbf{迁移信号：}给出 $2p$ 与角度，不必求交点。\par\medskip
\textbf{\textcolor{CExample}{例2（反求倾角，注意两解）}}\par
\textbf{题目：}$y^2=8x$ 的焦点弦长为 $16$，求所在直线的倾斜角。\par
\textbf{解：}$2p=8$，故 $\sin^2\theta=8/16=1/2$。
因 $0<\theta<\pi$，
\[
\boxed{\theta=\frac\pi4\ \text{或}\ \frac{3\pi}4.}
\]
\textbf{迁移信号：}正弦平方相同，角度通常成互补的一对。\par\medskip
\textbf{\textcolor{CExample}{例3（焦半径相加，逆求参数）}}\par
\textbf{题目：}$y^2=2px$，焦点弦 $AB$ 的倾角是 $60^\circ$，且
$|AF|+|BF|=8$，求 $p$。\par
\textbf{解：}焦点在两端点之间，$L=|AF|+|BF|=8$。因此
\[
8=\frac{2p}{\sin^260^\circ}=\frac{8p}{3}
\quad\Longrightarrow\quad\boxed{p=3}.
\]
\textbf{迁移信号：}遇到两段焦半径之和，立即替换为整条弦长。\par\medskip
\end{examplebox}
\clearpage
\begin{examplebox}
\textbf{\textcolor{CExample}{例4（左开口，公式不换号）}}\par
\textbf{题目：}$y^2=-6x$，焦点弦所在直线倾角 $120^\circ$，求 $L$。\par
\textbf{解：}$2p=6$，$\sin^2120^\circ=3/4$，
\[\boxed{L=\frac6{3/4}=8}.\]
\textbf{陷阱：}尽管焦点横坐标是负的，$p$ 一直取正数。\par\medskip
\textbf{\textcolor{CExample}{例5（给斜率时直接算；竖直另看）}}\par
\textbf{题目：}抛物线 $y^2=8x$，焦点弦所在直线斜率 $k=2$，求弦长；若直线竖直，弦长是多少？\par
\textbf{解：}由 $2p=8$ 得
\[
\boxed{L=\frac{8(1+2^2)}{2^2}=10}.
\]
竖直时 $\theta=\pi/2$，是通径，故 $\boxed{L=8}$。
\textbf{陷阱：}竖直直线没有有限斜率，不能往 $k$ 公式中硬代入。\par\medskip
\end{examplebox}
\clearpage
\begin{examplebox}
\textbf{\textcolor{CExample}{例6（坐标核验：检查图像是否真的准确）}}\par
\textbf{题目：}$y^2=4x$，$F(1,0)$，过 $F$ 且倾角 $60^\circ$ 的直线交抛物线于 $A,B$。给出交点并核对弦长。\par
\textbf{解：}直线可设 $x=1+y/\sqrt3$，代入 $y^2=4x$ 得
$y^2-\frac4{\sqrt3}y-4=0$，
其两根为 $2\sqrt3$ 与 $-\frac{2\sqrt3}{3}$，故
\[
A=(3,2\sqrt3),\qquad B=(\tfrac13,-\tfrac{2\sqrt3}{3}).
\]
两点满足 $y^2=4x$ 且在所设直线上，弦长
\[
\begin{aligned}
|AB|&=\sqrt{(3-\tfrac13)^2+(2\sqrt3+\tfrac{2\sqrt3}3)^2}\\
&=\sqrt{(\tfrac83)^2+(\tfrac{8\sqrt3}3)^2}=\boxed{\tfrac{16}3}
=\frac4{\sin^260^\circ}.
\end{aligned}
\]
\textbf{回看理解说明的总览图：}蓝段 $FA=4$，橙段 $FB=4/3$，与坐标距离所得 $16/3$ 完全一致。\par\medskip
\textbf{\textcolor{CExample}{高价值变式（检测是否真正理解）}}\par
\textbf{题目：}某过焦点弦长恰是通径的 $4$ 倍，问倾斜角及两侧焦半径之比（右开口）？\par
\textbf{解：}$L/(2p)=4$，所以 $\sin^2\theta=1/4$，
$\theta=30^\circ$ 或 $150^\circ$。
方法一中 $r:s=(1+\cos\theta):(1-\cos\theta)$，
于是 $\theta=30^\circ$ 时
\[
\boxed{r:s=(2+\sqrt3):(2-\sqrt3)},
\]
$150^\circ$ 时顺序交换。\textbf{检查：}两种角的整条弦长一样，两侧长短对调。
\end{examplebox}
